% =============================================================================
% Notes file: Hallucination operationalization for CLD generation
% This content has been integrated into ch2_background.tex, Section 2.7
% Kept here as reference notes
% =============================================================================

% EDGE-LEVEL FOCUS
% - Hallucinations operationalized at causal link (edge) level
% - Enables precise quantification
% - Aligns with atomic unit that judges can evaluate

% HALLUCINATION TAXONOMY (adapted from SD-AI benchmark)
% Reference: \cite{schoenberg2025aibuildsdmodels}
% Leaderboard: https://ub-iad.github.io/sd-ai/#/leaderboard/cld

% Types addressed in this thesis:
% 1. Spurious edge (FP) - fabricated causal relationship
% 2. Missing edge (FN) - omitted valid relationship  
% 3. Polarity error - wrong sign (+/-)
% 4. Wrong causal narrative - low judge score

% Types NOT primary focus:
% - Lack of conformance - solved via structured outputs (JSON schema)
% - Variable extraction failures - see Appendix D

% GROUND TRUTH DEFINITIONS
% 1. GT Lit: FP/FN w.r.t. literature CLD (peer-reviewed)
% 2. GT Synth: Corrupted edges = hallucinations (HaluEval approach)

% CORRUPTION PROCESS (GT Synth)
% Reference: \cite{li2023halueval}
% - Same model family as generator/judge (GPT-4.1) for self-preference bias
% - Corruption rate: 0.3 (30% of eligible edges)
% - Balanced distribution across types:
%   a) Motivation corruption: plausible-but-flawed reasoning
%   b) Spurious edge injection: NONE -> false positive
%   c) Polarity flip: POSITIVE <-> NEGATIVE

% CORRUPTOR PARAMETERS
% model=gpt-4.1, temperature=0.7, top_p=1.0, corruption_rate=0.3
